\section{Fine cells and sparse belts}
\label{sec:area_law_fine_belts}

The origin $o\in\mathbb R^2$ remains arbitrary. All scale indices and
refinement depths are nonnegative integers. Let $Z\subset\mathbb Z^2$ be
finite, and let $C_0\ge0$ be an integer. Retain the layer $D_k$ and its
index set $I_k$ from Definition~\ref{def:area_law_dyadic_layers}.

\begin{definition}[Dyadic refinement]
    \label{def:area_law_dyadic_refinement}
    \lean{TNLean.PEPS.AreaLaw.Geometry.dyadicAncestor,
        TNLean.PEPS.AreaLaw.Geometry.dyadicRefinement}
    \leanok
    For every refinement depth $d$ and fine index $u\in\mathbb Z^2$, define
    its coarse ancestor by
    \begin{align}
        A_d(u) & =\left(\left\lfloor\frac{u_1}{2^d}\right\rfloor,
        \left\lfloor\frac{u_2}{2^d}\right\rfloor\right).
    \end{align}
    For every finite set $P\subset\mathbb Z^2$, define its refinement by
    \begin{align}
        R_d(P) & =\{2^d z+a:z\in P,\ a\in\{0,\ldots,2^d-1\}^2\}.
    \end{align}
\end{definition}

\begin{theorem}[Ancestors and refinement counts]
    \label{thm:area_law_dyadic_refinement}
    \lean{TNLean.PEPS.AreaLaw.Geometry.mem_dyadicRefinement_iff,
        TNLean.PEPS.AreaLaw.Geometry.card_dyadicRefinement,
        TNLean.PEPS.AreaLaw.Geometry.dyadicCellIndex_of_le}
    \leanok
    \uses{def:area_law_dyadic_refinement,def:area_law_dyadic_cells}
    For every finite $P$, depth $d$, and index $u\in\mathbb Z^2$,
    \begin{align}
        u\in R_d(P) & \quad\Longleftrightarrow\quad A_d(u)\in P, \\
        |R_d(P)|    & =4^d|P|.
    \end{align}
    For every pair of scales $\ell\le k$ and every $x\in\mathbb R^2$,
    \begin{align}
        A_{k-\ell}(c_\ell(x)) & =c_k(x).
    \end{align}
\end{theorem}
\begin{proof}\leanok
    Euclidean division by $2^d$ gives a unique integer quotient and a
    remainder in $\{0,\ldots,2^d-1\}$ in each coordinate. This proves the
    membership criterion and the cardinality formula, including negative
    indices. The identity $\lfloor t/n\rfloor=\lfloor\lfloor t\rfloor/n\rfloor$
    for a positive integer $n$ gives the point-index identity.
\end{proof}

\begin{definition}[Fine indices of a layer]
    \label{def:area_law_fine_layer_indices}
    \lean{TNLean.PEPS.AreaLaw.Geometry.fineLayerIndices}
    \leanok
    \uses{def:area_law_dyadic_refinement,def:area_law_dyadic_layers}
    For every fine scale $\ell\le k$, set
    \begin{align}
        F_{k,\ell} & =R_{k-\ell}(I_k).
    \end{align}
\end{definition}

\begin{theorem}[Exact fine-cell layer union]
    \label{thm:area_law_fine_layer_count}
    \lean{TNLean.PEPS.AreaLaw.Geometry.dyadicLayer_eq_iUnion_fine,
        TNLean.PEPS.AreaLaw.Geometry.card_fineLayerIndices,
        TNLean.PEPS.AreaLaw.Geometry.card_fineLayerIndices_le,
        TNLean.PEPS.AreaLaw.Geometry.card_fineLayerIndices_boundary_le}
    \leanok
    \uses{def:area_law_fine_layer_indices,def:area_law_dyadic_cells}
    For every $\ell\le k$, the actual layer and its fine-cell count satisfy
    \begin{align}
        D_k          & =\bigcup_{u\in F_{k,\ell}} C_{\ell,u}, \\
        |F_{k,\ell}| & =4^{k-\ell}|I_k|,                      \\
        |F_{k,\ell}| & \le4(2C_0+1)^2|Z|\,4^{k-\ell}.
    \end{align}
    If $Z$ is the set of endpoints of the crossing edges of a cut in a
    finite induced domain and $b$ is the number of those edges, then
    \begin{align}
        |F_{k,\ell}| & \le8(2C_0+1)^2b\,4^{k-\ell}.
    \end{align}
    These formulas include empty layers and equal scales. The half-open
    fine cells are disjoint and nonempty, so $F_{k,\ell}$ is exactly the
    set of fine indices whose cells meet $D_k$. These are the predetermined
    cells used in the belt count of
    \cite[Section~11]{OpenAI2026AreaLaw}.
\end{theorem}
\begin{proof}\leanok
    \uses{thm:area_law_dyadic_refinement,thm:area_law_dyadic_indices,
        thm:area_law_dyadic_nesting,thm:area_law_dyadic_layer_count}
    A point's fine index belongs to $R_{k-\ell}(I_k)$ precisely when its
    coarse index belongs to $I_k$. The cell-membership criterion therefore
    proves the exact union. Multiply the coarse-cell count by $4^{k-\ell}$
    and apply the endpoint and crossing-edge bounds.
\end{proof}

\begin{definition}[Fine scale and belt pitch]
    \label{def:area_law_fine_belt_scales}
    \lean{TNLean.PEPS.AreaLaw.Geometry.fineScaleIndex,
        TNLean.PEPS.AreaLaw.Geometry.pitchScaleIndex}
    \leanok
    \uses{thm:area_law_exponent_gaps}
    With the fixed exponents $\delta_0$ and $\zeta=1-\delta_0/2$, define,
    for every scale $k$,
    \begin{align}
        \ell_k & =\lfloor\zeta k\rfloor,       \\
        p_k    & =\lfloor(1+\delta_0)k\rfloor, \\
        r_k    & =2^k,                         \\
        t_k    & =2^{\ell_k},                  \\
        s_k    & =2^{p_k}.
    \end{align}
    Both floor arguments are nonnegative, so the displayed floors are
    the ordinary integer floors. The fine cells have side $t_k$, while
    $s_k$ is the belt pitch, as in
    \cite[Section~11]{OpenAI2026AreaLaw}.
\end{definition}

\begin{theorem}[Scale order, divisibility and rounding]
    \label{thm:area_law_fine_belt_scales}
    \lean{TNLean.PEPS.AreaLaw.Geometry.fineScaleIndex_le,
        TNLean.PEPS.AreaLaw.Geometry.le_pitchScaleIndex,
        TNLean.PEPS.AreaLaw.Geometry.fineScale_dvd_layerScale,
        TNLean.PEPS.AreaLaw.Geometry.layerScale_dvd_pitchScale,
        TNLean.PEPS.AreaLaw.Geometry.layerScale_exponent_le}
    \leanok
    \uses{def:area_law_fine_belt_scales}
    For every scale $k$,
    \begin{align}
        \ell_k        & \le k\le p_k,              \\
        t_k           & \mid r_k,                  \\
        r_k           & \mid s_k,                  \\
        2k-p_k-\ell_k & \le2-\frac{\delta_0 k}{2}.
    \end{align}
\end{theorem}
\begin{proof}\leanok
    \uses{thm:area_law_exponent_gaps}
    The inequalities $\zeta\le1$ and $\delta_0\ge0$ give the scale order.
    Ordered exponents give divisibility of the corresponding powers of
    two. Each floor is greater than its argument minus one. Adding the
    two floor bounds and using $\zeta=1-\delta_0/2$ proves the last inequality.
\end{proof}

\begin{definition}[Belt cells]
    \label{def:area_law_belt_cells}
    \lean{TNLean.PEPS.AreaLaw.Geometry.beltCellIndices}
    \leanok
    Let $F\subset\mathbb Z^2$ be finite and let $m\ge1$ be an integer.
    For every pair of residues $a,b\in\{0,\ldots,m-1\}$, define
    \begin{align}
        B(F;m,a,b) & =\{u\in F:u_1\equiv a\pmod m\ \text{or}\ u_2\equiv b\pmod m\}.
    \end{align}
    For a fine layer with pitch exponent $p\ge k\ge\ell$, the modulus is
    $m=2^{p-\ell}$, the ratio of pitch to fine-cell side.
\end{definition}

\begin{theorem}[Sparse shift in the actual fine-cell layer]
    \label{thm:area_law_fine_belt_shift}
    \lean{TNLean.PEPS.AreaLaw.Geometry.exists_fine_belt_shift}
    \leanok
    \uses{def:area_law_belt_cells,def:area_law_fine_layer_indices}
    Let $Z$ be the crossing-edge endpoint set of a finite induced-domain
    cut, with crossing-edge count $b$. For every pair of scales
    $\ell\le k$ and pitch exponent $p\ge k$, put $F=F_{k,\ell}$ and
    $m=2^{p-\ell}$. There are residues $a,b'\in\{0,\ldots,m-1\}$ such that
    \begin{align}
        m\,|B(F;m,a,b')| & \le16(2C_0+1)^2b\,4^{k-\ell}.
    \end{align}
    Thus, with $r=2^k$, $t=2^\ell$, and $s=2^p$,
    \begin{align}
        |B(F;m,a,b')| & \le16(2C_0+1)^2b\,\frac{r^2}{st}.
    \end{align}
\end{theorem}
\begin{proof}\leanok
    \uses{thm:area_law_sparse_belt_shift,thm:area_law_fine_layer_count}
    Apply coordinate-residue averaging to the fixed set $F$. The selected
    belt cells satisfy $m|B|\le2|F|$. The actual fine-cell count gives the
    first inequality. Divide by the positive modulus and use
    $4^{k-\ell}/2^{p-\ell}=r^2/(st)$ for the second.
\end{proof}

\begin{theorem}[Geometric decay of the scale ratio]
    \label{thm:area_law_fine_belt_ratio}
    \lean{TNLean.PEPS.AreaLaw.Geometry.dyadicScale_ratio_le}
    \leanok
    \uses{def:area_law_fine_belt_scales}
    For every scale $k$,
    \begin{align}
        \frac{r_k^2}{s_k t_k} & \le4\,2^{-\delta_0 k/2}.
    \end{align}
\end{theorem}
\begin{proof}\leanok
    \uses{thm:area_law_fine_belt_scales}
    The ratio equals $2^{2k-p_k-\ell_k}$. The exponent bound gives
    $2^{2k-p_k-\ell_k}\le2^{2-\delta_0 k/2}=4\,2^{-\delta_0 k/2}$.
    The factor four accounts for the two floor-rounding losses.
\end{proof}

% Source: geometry:belt-count in Section 11.
\begin{theorem}[Sparse dyadic belts]
    \label{thm:area_law_sparse_dyadic_belts}
    \lean{TNLean.PEPS.AreaLaw.Geometry.exists_sparse_dyadic_belt_shift}
    \leanok
    \uses{def:area_law_belt_cells,def:area_law_fine_layer_indices,
        def:area_law_fine_belt_scales}
    Let $Z$ be the crossing-edge endpoint set of a finite induced-domain
    cut, with crossing-edge count $b$. For every scale $k$, put
    $F=F_{k,\ell_k}$ and $m=2^{p_k-\ell_k}=s_k/t_k$. There are residues
    $a,b'\in\{0,\ldots,m-1\}$ such that
    \begin{align}
        |B(F;m,a,b')| & \le64(2C_0+1)^2b\,2^{-\delta_0 k/2}.
    \end{align}
    This is the sparse-belt estimate in \cite[Section~11]{OpenAI2026AreaLaw}.
    The constant is fixed after the radius $C_0$ and is independent of the
    domain, cut, origin and scale. The residues are chosen separately in
    each layer. Empty endpoint sets and layers are included.
\end{theorem}
\begin{proof}\leanok
    \uses{thm:area_law_fine_belt_shift,thm:area_law_fine_belt_scales,
        thm:area_law_fine_belt_ratio}
    The scale order permits the fine-cell shift theorem with
    $\ell=\ell_k$ and $p=p_k$. Its coefficient $16(2C_0+1)^2$ multiplied
    by the scale-ratio factor four gives $64(2C_0+1)^2$.
\end{proof}
